\section{问题三模型的建立与求解}

\subsection{问题三的描述与分析}

问题三固定计算任务排程，要求为六个区域设计储能充放电策略，并分析储能对运行成本、碳排放、区域峰值净购电功率和负荷波动的影响。根据题意，设施负荷不再是决策变量，而由附件给定的 AI 与 NonAI IT 负荷直接确定：
\begin{equation}
L[r,t]=\mathrm{PUE}[r]\Bigl(\mathrm{BaselineAI\_IT\_Load\_MW}[r,t]+\mathrm{NonAI\_IT\_Load\_MW}[r,t]\Bigr),
\label{eq:q3-fixed-load}
\end{equation}
其中 $r\in\{\mathrm{A},\ldots,\mathrm{F}\}$，$t=0,1,\ldots,2406$。本问只优化新能源直接供负荷、储能充放电、电网购电、新能源售电与弃电。

由于设施负荷被固定，储能不可能改变计算负荷本身的波动，因此题目中的“负荷波动”需要区分两层含义：\textbf{设施负荷波动在优化前后严格不变；储能能够改变的是数据中心与外部电网之间的净交互功率波动。} 后文同时报告二者，避免把电网侧平滑效果误写成计算负荷下降。

本问先从附件数据识别区域储能角色，而不预设六区采用同一种策略。六区总可用新能源约为总设施负荷的 1.74--1.81 倍，A--E 全时段新能源均不低于固定负荷，F 仅有 1 个短缺小时；A/B/C 不允许售电，而 D/E/F 允许外送。由此，储能价值主要取决于新能源盈余、售电边界和局部短缺，而不是区域名称。总体流程见图\ref{fig:q3-flow}。

\begin{figure}[!htbp]
\centering
\includegraphics[width=\FigureFlowWidth]{\ProjectRoot/modules/40_q3/figures/q3_flow.png}
\caption{问题三从能源结构诊断到储能经济调度的总体流程}
\label{fig:q3-flow}
\end{figure}

\subsection{数据结构诊断与储能角色识别}

\subsubsection{附件能源守恒与基准可比性}

对附件 $6\times2407$ 个区域—小时逐点检查，新能源分配、设施负荷、充电分解、净电网交互和碳排放恒等式均可在约 $10^{-4}$--$2\times10^{-4}$ 的数值误差内复现。因此，正式模型沿用附件的新能源分配、负荷平衡、电网购电分解和 SOC 动态，不另造能源流语义。

附件原始运行状态存在一个重要可比性问题：D/E/F 的 $SOC(2406)$ 低于 InitialSOC，RegionE 在 Hour~0 还存在约 1 MWh 的单点 SOC 不一致。因此，\textbf{附件原始能源基准只能作为参考状态，不能直接视为满足问题三终端约束的可行优化点。} 为区分“当期新能源重新分配”与“储能跨时搬移”两类收益，必须再设置一个同口径的无储能优化对照。

\subsubsection{新能源盈余结构}

定义固定设施负荷下的新能源盈余
\begin{equation}
H[r,t]=\mathrm{AvailableRenewable}[r,t]-L[r,t].
\label{eq:q3-surplus}
\end{equation}
表\ref{tab:q3-surplus-structure}给出六区域的核心结构。

\begin{table}[!htbp]
\centering
\caption{六区域固定负荷下的新能源盈余结构}
\label{tab:q3-surplus-structure}
\begin{tabular}{lccc}
\toprule
\rowcolor{CumcmLightBlue}
区域 & 总可用新能源/总设施负荷 & $H<0$ 小时数 & 售电能力 \\
\midrule
A & 1.756 & 0 & 不允许售电 \\
B & 1.775 & 0 & 不允许售电 \\
C & 1.745 & 0 & 不允许售电 \\
D & 1.770 & 0 & 允许售电 \\
E & 1.812 & 0 & 允许售电 \\
F & 1.782 & 1 & 允许售电 \\
\bottomrule
\end{tabular}
\end{table}
\FloatBarrier

A--E 的 2407 个小时均满足 $H[r,t]\ge0$；RegionF 仅 1 个小时 $H[r,t]<0$。因此，多数区域并不存在能源总量不足，真正决定储能价值的是售电边界、售电价格时序以及局部新能源短缺。

\subsubsection{三类区域储能角色}

由 $H[r,t]$ 的符号和 SellLimit 可直接得到三类区域：
\begin{enumerate}
  \item \textbf{A/B/C：持续盈余且不可售电。} 新能源始终可覆盖负荷，多余电量只能弃置；在成本最优且储能动作最少的代表解中，无需储能即可实现零购电，因此应出现零充放电动作。
  \item \textbf{D/E：持续盈余且允许售电。} 储能主要承担新能源价值时移：低售电价值时段用新能源充电，高售电价值时段由储能承担部分设施负荷，从而释放同期新能源用于外送。
  \item \textbf{F：盈余为主、存在单次短缺且允许售电。} 除售电价值时移外，还可在唯一短缺小时通过放电替代电网购电。
\end{enumerate}

上述分类由求解前的数据结构得到，同一线性规划在六区保持完全一致的数学形式；分类仅用于解释为什么统一模型会产生不同的最优行为。

\subsubsection{价格、碳强度与新能源的时序关系}

在\textbf{单一区域内部沿时间维度}计算，0--2399\,h 的 Price 与 CarbonIntensity 相关系数约为 $-0.378$，CarbonIntensity 与 AvailableRenewable 的相关约为 $-0.663$。这与问题二将六区域—小时\textbf{整体混合统计}后得到的 Price--Carbon 正相关并不矛盾：问题二的 pooled 相关同时包含区域间平均价格和平均碳水平差异，而本节只考察各区域内部的时间变化。

如果系统仍需大量购电，低价时段并不一定低碳；但本题固定负荷下，无储能优化后系统购电已降至仅 F 区约 31.60 MWh，加入储能后进一步降为数值零。因此 Carbon 在问题三中主要随购电量一起退化为结果指标，没有必要为成本与碳排构造人工权重。

综合上述数据，问题三的模型结构可概括为：固定设施负荷下先识别 $H$ 与售电边界，再用统一能源流线性规划完成当期分配和储能跨时搬移；为避免纯成本 LP 的多重最优造成无意义储能循环，采用“\textbf{净运行成本最小} $\rightarrow$ \textbf{储能吞吐量最小}”两层字典序选择物理代表解。

\subsection{固定负荷下的能源—储能调度模型}

\subsubsection{决策变量}

对区域 $r$、小时 $t$ 定义非负连续变量：$u[r,t]$ 为新能源直接供负荷功率，$qR[r,t]$ 为新能源充电功率，$qG[r,t]$ 为电网充电功率，$d[r,t]$ 为储能放电功率，$gL[r,t]$ 为电网直接供负荷功率，$s[r,t]$ 为新能源售电功率，$w[r,t]$ 为弃电功率，$E[r,t]$ 为小时末 SOC。

派生量为
\begin{equation}
\mathrm{GridPurchase}=gL+qG,\qquad
\mathrm{ChargePower}=qR+qG,\qquad
\mathrm{NetGridImport}=\mathrm{GridPurchase}-s.
\label{eq:q3-derived}
\end{equation}

\subsubsection{能量守恒与 SOC 动态}

新能源分配满足
\begin{equation}
\mathrm{AvailableRenewable}[r,t]=u[r,t]+qR[r,t]+s[r,t]+w[r,t],
\label{eq:q3-renewable-balance}
\end{equation}
固定设施负荷满足
\begin{equation}
L[r,t]=u[r,t]+d[r,t]+gL[r,t],
\label{eq:q3-load-balance}
\end{equation}
且
\begin{equation}
\mathrm{GridPurchase}[r,t]=gL[r,t]+qG[r,t].
\label{eq:q3-grid-split}
\end{equation}
售电 $s[r,t]$ 仅允许来自当期新能源分配，不允许网购电或储能放电直接售电。

SOC 动态为
\begin{equation}
E[r,t]=E[r,t-1]+\eta_c[r]\bigl(qR[r,t]+qG[r,t]\bigr)\Delta t-\frac{d[r,t]}{\eta_d[r]}\Delta t,
\label{eq:q3-soc-dynamics}
\end{equation}
其中
\begin{equation}
E[r,-1]=\mathrm{InitialSOC}[r],\qquad E[r,2406]\ge\mathrm{InitialSOC}[r].
\label{eq:q3-terminal-soc}
\end{equation}
储能和电网边界为
\begin{align}
&\mathrm{MinSOC}[r]\le E[r,t]\le\mathrm{StorageCapacity}[r],\label{eq:q3-soc-bound}\\
&0\le qR[r,t]+qG[r,t]\le\mathrm{MaxChargePower}[r],\label{eq:q3-charge-bound}\\
&0\le d[r,t]\le\mathrm{MaxDischargePower}[r],\label{eq:q3-discharge-bound}\\
&0\le\mathrm{GridPurchase}[r,t]\le\mathrm{MaxGridImport}[r],\label{eq:q3-grid-import}\\
&0\le s[r,t]\le\min\bigl(\mathrm{SellLimit}[r],\mathrm{MaxGridExport}[r]\bigr).\label{eq:q3-sell-bound}
\end{align}

\subsubsection{两层字典序目标}

第一层最小化净运行成本：
\begin{equation}
\min\ \mathrm{Cost}=\sum_{r,t}\Bigl[\mathrm{Price}[r,t]\bigl(gL[r,t]+qG[r,t]\bigr)-\mathrm{SellPrice}[r,t]s[r,t]\Bigr]\Delta t.
\label{eq:q3-stage1-objective}
\end{equation}
售电收益可能大于购电支出，因此净运行成本允许为负，负值表示净售电收益。

纯成本 LP 可能存在多重最优解。为从成本等价解中选择储能动作最少、物理意义最清楚的方案，第二层在成本最优面上最小化储能总吞吐量：
\begin{equation}
\min\ \mathrm{Throughput}=\sum_{r,t}\bigl[qR[r,t]+qG[r,t]+d[r,t]\bigr]\Delta t,
\label{eq:q3-stage2-objective}
\end{equation}
并增加
\begin{equation}
\mathrm{Cost}\le\mathrm{Cost}^{*}+10^{-4}\ \mathrm{CNY}.
\label{eq:q3-stage2-cost-tolerance}
\end{equation}
这样无需人为加入电池退化成本或二元变量，也能消除成本等价的同时充放电循环。

\subsubsection{模型汇总}

问题三的统一模型可写为
\begin{equation}
\left\{
\begin{aligned}
&\text{第一层：}\min\ \sum_{r,t}[\mathrm{Price}(gL+qG)-\mathrm{SellPrice}\,s]\Delta t,\\
&\text{第二层：}\min\ \sum_{r,t}(qR+qG+d)\Delta t,\\
&\mathrm{AvailableRenewable}=u+qR+s+w,\\
&L=u+d+gL,\qquad \mathrm{GridPurchase}=gL+qG,\\
&E[t]=E[t-1]+\eta_c(qR+qG)-d/\eta_d,\\
&\mathrm{MinSOC}\le E[t]\le\mathrm{StorageCapacity},\\
&0\le qR+qG\le\mathrm{MaxChargePower},\quad0\le d\le\mathrm{MaxDischargePower},\\
&0\le\mathrm{GridPurchase}\le\mathrm{MaxGridImport},\\
&0\le s\le\min(\mathrm{SellLimit},\mathrm{MaxGridExport}),\\
&E[2406]\ge\mathrm{InitialSOC},\qquad\text{其余能源流变量非负}.
\end{aligned}
\right.
\label{eq:q3-model-summary}
\end{equation}
模型中不存在跨区域能量交换，因此六个区域可严格分解为 6 个独立线性规划并行求解，再汇总系统指标。

\subsection{模型求解与三层对照}

为准确识别不同机制的贡献，设置三种物理意义明确的方案：
\begin{enumerate}
  \item \textbf{附件原始能源基准}：直接使用附件给定的运行状态，仅作为现实参考；其 D/E/F 终端 SOC 不满足本问统一终端约束。
  \item \textbf{固定负荷无储能优化方案}：保持相同负荷和能源参数，令 $qR=qG=d=0$，仅优化新能源直接利用、售电、弃电和电网购电。该方案衡量“重新分配当期新能源”本身的价值。
  \item \textbf{固定负荷储能协同方案}：使用完整模型，允许储能跨时段搬移新能源。
\end{enumerate}
为与程序和结果文件对应，上述三层分别对应内部结果标识 B0、E0、E1；\textbf{正文后续仅使用其实际物理名称，不再以代号作为模型名称。}

储能的独立增量价值严格定义为“储能协同方案 $-$ 无储能优化方案”的指标变化；“附件原始基准 $-$ 无储能优化方案”反映原始运行策略与同口径能源再优化之间的差距，不归因于储能。

每个区域独立求解上述两层线性规划，具体流程见算法\ref{alg:q3-two-stage}。

固定负荷下的储能调度可写成连续能源流优化，并通过购售电价格与荷电状态约束刻画跨时价值\cite{zhang_dc_bess_2025}。对于线性成本模型中可能出现的同时充放电现象，可通过次级目标消除无意义循环\cite{wang_exact_relax_2024}。因此采用“先最小成本、再最小吞吐量”的两层字典序求解。

\begin{algorithm}[H]
\SetAlgoLined
\setstretch{0.86}
\SetAlFnt{\small}
\SetArgSty{textnormal}
\SetKwInput{KwIn}{输入}
\SetKwInput{KwInit}{初始化}
\SetKwInput{KwOut}{输出}
\SetKwFor{ForEach}{对每个}{执行}{结束}
\caption{固定负荷下的两层字典序储能经济调度}
\label{alg:q3-two-stage}
\KwIn{设施负荷 $L_{rt}$，可用新能源 $R_{rt}$，购售电价格、碳强度，储能与电网参数}
\KwInit{$\varepsilon\leftarrow10^{-4}$ CNY；设置附件原始方案、无储能优化方案和 BESS-ED 方案}

\ForEach{区域 $r\in\{A,B,C,D,E,F\}$}{
  读取 $L_{rt}$、$R_{rt}$、价格及储能参数，构造 $u,qR,qG,d,gL,s,w,E$ 的变量索引\;
  建立新能源平衡、负荷平衡和 SOC 递推等式，并加入充放电功率、购售电上限及终端 SOC 约束\;

  \textbf{第一层：最小化净运行成本}\;
  构造
  $\mathrm{Cost}_r=\sum_t[c^{buy}_{rt}(gL_{rt}+qG_{rt})-c^{sell}_{rt}s_{rt}]\Delta t$\;
  求解区域线性规划；若不可行则返回约束信息，否则记录最优值 $\mathrm{Cost}_r^*$\;

  \textbf{第二层：在成本最优集合内最小化储能吞吐量}\;
  增加 $\mathrm{Cost}_r\le\mathrm{Cost}_r^*+\varepsilon$\;
  最小化 $\mathrm{Throughput}_r=\sum_t(qR_{rt}+qG_{rt}+d_{rt})\Delta t$\;

  重构 $G^{buy}_{rt}=gL_{rt}+qG_{rt}$、$P^{net}_{rt}=G^{buy}_{rt}-s_{rt}$ 与逐时 SOC\;
  计算区域成本、碳排放、新能源利用率、峰值净购电、净电网波动、平均绝对爬坡及储能吞吐量\;
  检查新能源/负荷平衡、SOC 递推、终端 SOC、充放电功率、购售电边界与同时充放电小时数\;
}

汇总六区域结果，并以“储能协同方案减无储能优化方案”的差值计算储能增量价值\;
\KwOut{六区域逐时储能与购售电方案、系统评价指标及约束检验结果}
\end{algorithm}


\subsection{结果与分析}

\subsubsection{系统总体结果}

表\ref{tab:q3-system-results}给出三种方案的系统级结果。

\begin{table}[!htbp]
\centering
\caption{固定负荷下三种能源调度方案的系统级结果}
\label{tab:q3-system-results}
\begin{tabular}{lrrr}
\toprule
\rowcolor{CumcmLightBlue}
指标 & 附件原始能源基准 & 无储能优化方案 & 储能协同方案 \\
\midrule
净运行成本/CNY & 1,801,660,137.89 & $-419,820,327.74$ & $-459,217,791.04$ \\
碳排放/tCO$_2$ & 2,045,367.48 & 8.8506 & 数值零 \\
新能源利用率 & 32.86\% & 67.91\% & 69.51\% \\
电网购电/MWh & 4,134,366.12 & 31.60 & 数值零 \\
新能源售电/MWh & 1,375,256.18 & 1,330,267.77 & 1,491,718.26 \\
电网交互功率标准差/MW & 607.20 & 115.39 & 4.98 \\
电网交互平均绝对爬坡/MW & 153.78 & 29.29 & 0.40 \\
\bottomrule
\end{tabular}
\end{table}
\FloatBarrier

储能协同方案相对无储能优化方案进一步改善净运行成本 39,397,463.30 CNY，这才是储能的独立增量经济价值。附件原始基准到无储能方案已经出现数量级更大的成本和碳排改善，说明附件新能源总量非常充足，原始运行策略差距不能全部算作储能收益。

碳排严格由电网购电量与碳强度计算。无储能方案仅 RegionF 购电 31.60 MWh，对应 8.8506 tCO$_2$；储能协同方案六区购电均降为数值零，因此碳排也为数值零。

区域峰值净购电定义为
\begin{equation}
\mathrm{PeakImport}[r]=\max_t\max\bigl(\mathrm{NetGridImport}[r,t],0\bigr).
\label{eq:q3-peak-import}
\end{equation}
无储能方案的最大区域峰值净购电发生在 RegionF，为 31.60 MW；储能协同方案六区均为数值零。系统汇总净交换可能因不同区域购售电抵消，因此不能替代区域峰值净购电指标。

\subsubsection{电网交互波动}

设施负荷由附件固定，其系统标准差在三种方案下均为 563.90 MW，优化前后变化严格为 0。储能改变的是外部电网交互功率：从无储能方案到储能协同方案，系统净电网交互功率标准差由 115.39 MW 降至 4.98 MW，下降约 95.68\%；平均绝对爬坡由 29.29 MW 降至 0.40 MW，下降约 98.62\%。

\begin{figure}[!htbp]
\centering
\includegraphics[width=\FigureDoubleWidth]{\ProjectRoot/modules/40_q3/figures/q3_netgrid_daily.png}
\caption{无储能与储能协同方案的系统净电网交互功率日统计}
\label{fig:q3-netgrid-daily}
\end{figure}
\FloatBarrier

\subsubsection{分区储能角色验证}

A/B/C 在储能协同方案中选择零储能动作，购电和售电均为 0，与“持续盈余且不可售电”的数据结构一致。D/E 的储能主要进行新能源价值时移：低售电价值时段以新能源充电，高售电价值时段通过储能放电承担设施负荷，从而释放同期新能源外送；两区储能增量价值分别为 9,743,701.13 CNY 和 13,696,192.16 CNY。RegionF 除售电价值时移外，还在唯一短缺小时以储能放电替代购电，储能增量价值为 15,957,570.01 CNY。

\begin{figure}[!htbp]
\centering
\includegraphics[width=\FigureDoubleWidth]{\ProjectRoot/modules/40_q3/figures/q3_soc_daily_def.png}
\caption{D/E/F 区域储能荷电状态日统计}
\label{fig:q3-soc-daily}
\end{figure}
\FloatBarrier

\begin{table}[!htbp]
\centering
\caption{储能协同方案的区域储能运行指标}
\label{tab:q3-regional-storage}
\begin{tabular}{lcccc}
\toprule
\rowcolor{CumcmLightBlue}
区域 & 储能增量价值/CNY & 吞吐量/MWh & 充电量/MWh & 放电量/MWh \\
\midrule
A & 0.00 & 0.00 & 0.00 & 0.00 \\
B & 0.00 & 0.00 & 0.00 & 0.00 \\
C & 0.00 & 0.00 & 0.00 & 0.00 \\
D & 9,743,701.13 & 81,550.42 & 43,512.12 & 38,038.30 \\
E & 13,696,192.16 & 125,593.16 & 67,011.61 & 58,581.55 \\
F & 15,957,570.01 & 139,058.35 & 74,196.11 & 64,862.24 \\
\midrule
系统合计 & 39,397,463.30 & 346,201.93 & 184,719.85 & 161,482.09 \\
\bottomrule
\end{tabular}
\end{table}
\FloatBarrier

\begin{figure}[!htbp]
\centering
\includegraphics[width=\FigureMechanismWidth]{\ProjectRoot/modules/40_q3/figures/q3_f_mechanism.png}
\caption{RegionF 唯一新能源短缺小时的储能响应机制}
\label{fig:q3-f-mechanism}
\end{figure}
\FloatBarrier

第二层最小吞吐量优化后，在成本变化不超过 $10^{-4}$ CNY 的条件下，六区同时充放电区域—小时数稳定为 0，系统总储能吞吐量为 346,201.93 MWh。A/B/C 的微小数值购电与成本属于求解容差量级，正式展示统一记为 0。

\subsection{模型检验}

\subsubsection{约束满足性审计}

新能源平衡、负荷平衡与 SOC 递推的逐时最大残差均小于 $2\times10^{-13}$；终端 SOC、最小/最大 SOC、充放电功率、购电和售电边界均满足。第二层求解后同时充放电小时数为 0，说明字典序目标有效消除了纯成本最优面上的退化循环。

\subsubsection{求解一致性与机制验证}

完整 LP 与前期 reduced-model 探针在逐区储能增量价值上的差异小于 $10^{-4}$ CNY。更重要的是，求解前由 $H[r,t]$ 与 SellLimit 推出的 A/B/C、D/E、F 三类区域，分别对应正式结果中的零储能动作、售电价值时移、短缺覆盖三种行为，说明区域机制来自数据结构而不是事后命名。

本问结论依赖附件给定的新能源、负荷、电价和储能参数，不作一般性外推；但“先识别新能源盈余与售电边界，再用统一能源流模型验证区域机制”的建模路径可用于同类固定负荷储能调度问题。